% 附录 A 电磁学中的单位（Units in electromagnetism）
\chapter{电磁学中的单位}

电磁学中的单位是一片潜在的地雷阵。本书通篇使用 SI 单位制，但二十世纪的许多文献
使用非 SI 记号。因此本附录试图解开旧单位制的若干谜团，供在新单位制下长大、需要
解读较老文献的读者使用。

从十九世纪末到第二次世界大战结束后不久，cgs 单位制是关心此类事务的国际委员会
偏爱的单位制。在 cgs 制中，距离以厘米量度、质量以克、时间以秒——cgs 之名由此
而来。SI 制中我们用米（100 倍于厘米）、千克（1000 倍于克）与秒。距离与质量单位
的这一选择差异看似无关痛痒，最终却在其他单位中引起巨大变化。例如，力的量纲是
质量×长度/时间$^{2}$，故 cgs 制中力的单位达因比 SI 力的单位牛顿小
$1000\times100=10^{5}$ 倍。cgs 的能量单位是尔格（1 达因·厘米），比 SI 能量单位
焦耳小 $10^{5}\times10^{2}=10^{7}$ 倍。

不幸的是事情还没完：一谈到磁场的定义，情况就急剧变糟。原因在于两个单位制对场
B、H、M 的定义选择不同：
\begin{equation*}
\mathbf{B} = \mu_0(\mathbf{H}+\mathbf{M}) \qquad\text{（SI）}
\end{equation*}
\begin{equation*}
\mathbf{B} = \mathbf{H} + 4\pi\mathbf{M} \qquad\text{（cgs）}\tag{A.1}
\end{equation*}
于是在 cgs 制中 B 与 H 量纲相同，但单位却令人困惑地不同（前者以高斯（G）计，后者
以奥斯特（Oe）计）。H 与 M 关系的差异源于两种单位制对“距磁单极 q 为 r 处磁场
强度 H 的大小”的定义方程作了不同选择：
\begin{equation*}
H = \frac{q}{4\pi r^{2}} \quad\text{（SI）}\qquad H = \frac{q}{r^{2}} \quad\text{（cgs）}\tag{A.2}
\end{equation*}
省去 $4\pi$ 因子使这一方程以及磁偶极产生磁场的类似方程在 cgs 制中相当简洁。然而
$4\pi$ 因子并未从理论中完全消失，而是冒了出来，出现在式 A.1 中一个更不受欢迎的
地方。cgs 制按式 A.2 用并不存在的磁单极定义 H：一个单位磁单极在 1 厘米远处产生
1 奥斯特的场（若那里恰好潜伏着另一个单位磁单极，它将受到 1 达因的力）。SI 制中
H 的定义更有用——用电流来定义（因为人们在实验室里通常用电流源产生磁场，而不是
费心安排单位磁单极）。于是每米 n 匝、载流 I 安培的长螺线管在 SI 制中产生的场为
$H=nI$。

SI 制的逻辑是：让 $4\pi$ 出现在式 A.2 这类方程中，提醒我们点磁单极的场具有球对称
性（$4\pi$ 与球的表面积 $4\pi r^{2}$ 有关）。距点电荷 r 处的电场 E 也是如此：
\begin{equation*}
E = \frac{q}{4\pi\epsilon_0 r^{2}} \quad\text{（SI）}\qquad E = \frac{q}{r^{2}} \quad\text{（cgs）},\tag{A.3}
\end{equation*}
其中适用类似的论证。采用 SI，就必须忍受几个由对称性论证自然产生的额外因子出现在
点电荷与偶极子的场定义方程中；其余方程则干净得多、人为因子更少。

cgs 制唯一可说的好处是电场与磁场具有相同的量纲，某些相对论性方程因此看起来更
顺眼；且自由空间中 B=H 也有某种对称性（尽管这会带来混淆的机会）。但 cgs 制中
额外的 c 因子撒得到处都是：例如，电荷在电磁场中以速度 v 运动时的洛伦兹力方程
变为
\begin{equation*}
\mathbf{F} = q(\boldsymbol{\varepsilon} + \mathbf{v}\times\mathbf{B}) \quad\text{（SI）}
\end{equation*}
\begin{equation*}
\mathbf{F} = q\left(\boldsymbol{\varepsilon} + \frac{\mathbf{v}\times\mathbf{B}}{c}\right) \quad\text{（cgs）}.\tag{A.4}
\end{equation*}

更糟糕的还在后面。事实上有两种在 cgs 中建立电磁学的途径：静电单位（esu）或电磁
单位（emu），两者之间的换算需要转换因子。磁学中通常采用 emu 制。emu 的磁矩单位
通常就缩写为“emu”，是一个量纲为体积的单位（cgs 中为 cm$^{3}$）。磁化率
$\chi=M/H$ 在两种单位制中都无量纲，但在 cgs 中常写作 emu\,cm$^{-3}$——即使乍看
不像，它仍然是无量纲的。在 cgs 磁化率与 SI 磁化率之间换算时总有一个恼人的
$4\pi$ 因子要记住——它源于 cgs 制中的 $\mathbf{B}=\mathbf{H}+4\pi\mathbf{M}$
（见上）。例如退磁因子 N 无量纲，在 SI 中 $0<N<1$，在 cgs 中 $0<N<4\pi$。
遗憾的是 emu 在磁学研究的某些领域仍被广泛使用。表 A.1 总结了 SI 与 cgs 之间
相互换算最有用的结果。

玻尔磁子虽不是 SI 单位，$\mu_{\mathrm{B}}=9.274\times10^{-24}$ J\,T$^{-1}$ 却是
磁矩的有用量度——它对应氢原子 1s 电子的磁矩。对顺磁体，摩尔磁化率 $\chi_{\mathrm{m}}$
由居里定律给出（把式 2.44 稍作改写），SI 单位下为
\begin{equation*}
\chi_{\mathrm{m}} = \frac{\mu_0N_{\mathrm{A}}\mu_{\mathrm{eff}}^{2}\mu_{\mathrm{B}}^{2}}{3k_{\mathrm{B}}T}\tag{A.5}
\end{equation*}
$N_{\mathrm{A}}$ 是阿伏伽德罗常数。于是 $\chi_m T$ 不依赖温度（见图 2.10），可以把
它与有效磁矩联系起来。整理式 A.5 得
\begin{equation*}
\mu_{\mathrm{eff}} = \left[3k_{\mathrm{B}}/\mu_0N_{\mathrm{A}}\mu_{\mathrm{B}}^{2}\right]^{1/2}\sqrt{\chi_{\mathrm{m}}T},\tag{A.6}
\end{equation*}

\begin{table}[!ht]
\centering
\caption{SI 单位制与 cgs 单位制中的单位。缩写：m=米，g=克，N=牛顿，J=焦耳，
T=特斯拉，G=高斯，A=安培，Oe=奥斯特，Wb=韦伯，Mx=麦克斯韦。emu 是电磁单位
（electromagnetic unit）的缩写。注意磁化率在 SI 单位中无量纲。}
\begin{small}
\begin{tabular}{llcll}
\toprule
量 & 符号 & \multicolumn{3}{c}{换算} \\
\midrule
长度 & $x$ & $10^{-2}$ m & $=1$ & cm\\
质量 & $m$ & $10^{-3}$ kg & $=1$ & g\\
力 & $F$ & $10^{-5}$ N & $=1$ & dyne\\
能量 & $E$ & $10^{-7}$ J & $=1$ & erg\\
磁感应强度 & $\mathbf{B}$ & $10^{-4}$ T & $=1$ & G\\
磁场强度 & $\mathbf{H}$ & $10^{3}/4\pi$ A\,m$^{-1}$ & $=1$ & Oe\\
磁矩 & $\mu$ & $10^{-3}$ J\,T$^{-1}$ 或 A\,m$^{2}$ & $=1$ & erg\,G$^{-1}$\\
磁化强度（单位体积的矩） & $\mathbf{M}$ & $10^{3}$ A\,m$^{-1}$ 或 J\,T$^{-1}$\,m$^{-3}$ & $=1$ & Oe\\
磁化率 & $\chi$ & $4\pi$ & $=1$ & emu\,cm$^{-3}$\\
摩尔磁化率 & $\chi_m$ & $4\pi\times10^{-6}$ m$^{3}$\,mol$^{-1}$ & $=1$ & emu\,mol$^{-1}$\\
质量磁化率 & $\chi_g$ & $4\pi\times10^{-3}$ m$^{3}$\,kg$^{-1}$ & $=1$ & emu\,g$^{-1}$\\
磁通量 & $\phi$ & $10^{-8}$ T\,m$^{2}$ 或 Wb & $=1$ & G\,cm$^{-2}$ 或 Mx\\
退磁因子 & N & \multicolumn{2}{c}{$0<N<1$} & $0<N<4\pi$\\
\bottomrule
\end{tabular}
\end{small}
\end{table}

从而
\begin{equation*}
\mu_{\mathrm{eff}} = 797.8\sqrt{\chi_{\mathrm{m}}^{\mathrm{SI}}T}\tag{A.7}
\end{equation*}
\begin{equation*}
\mu_{\mathrm{eff}} = 2.827\sqrt{\chi_{\mathrm{m}}^{\mathrm{cgs}}T}\tag{A.8}
\end{equation*}
其中 $\mu_{\mathrm{eff}}$ 以每化学式单元玻尔磁子计，$\chi_{\mathrm{m}}^{\mathrm{SI}}$
以 m$^{3}$\,mol$^{-1}$ 计，$\chi_{\mathrm{m}}^{\mathrm{cgs}}$ 以 emu\,mol$^{-1}$
计。这些数值关系对从 $\chi_{\mathrm{m}}T$ 对 T 的图中提取有效磁矩很有用。

\section*{延伸阅读}

\begin{itemize}[itemsep=0.2em]
\item B.~I.~Bleaney 与 B.~Bleaney,《Electricity and Magnetism》, OUP 1989。
\item J.~D.~Jackson,《Classical Electrodynamics》, Wiley 1962。
\end{itemize}
